\documentclass[11pt,a4paper]{article}
\usepackage[margin=2.5cm]{geometry}
\usepackage[hyphens]{url}
\usepackage[hidelinks]{hyperref}
\usepackage[T1]{fontenc}
\usepackage[p,osf,swashQ]{cochineal}
\usepackage{enumitem}

\setlength{\parindent}{0pt}
\setlength{\parskip}{0.6em}

\begin{document}

\begin{flushright}
Ziliang Xiong\\
ziliang.xiong@liu.se \quad|\quad linkedin.com/in/xzleo \quad|\quad xzleo.github.io\\
Linköping University, Sweden
\end{flushright}

\vspace{1em}
\today

\vspace{1em}
Apple Machine Learning Research\\
Paris, France

\vspace{1em}
\textbf{Re: Internship -- Machine Learning Research (200683768)}

\vspace{0.5em}
Dear Hiring Team,

I am writing to apply for the Machine Learning Research internship at Apple MLR in Paris. I am a Ph.D. candidate in computer vision at Linköping University, affiliated with WASP and ELLIIT, and my research centers on \textbf{uncertainty quantification}---a topic explicitly listed in your posting. I am motivated by curiosity-driven, open research that advances both scientific understanding and trustworthy machine learning systems.

\textbf{Why this role fits my research.}
Your team works on foundational ML problems including uncertainty quantification, generative modeling, efficient inference, and data-centric learning. My Ph.D. applies UQ to computer vision tasks---regression, human pose estimation, and end-to-end autonomous driving---with the goal of making learned models safer and better calibrated. Representative work includes Hinge-Wasserstein for multimodal aleatoric uncertainty (CVPR 2024 workshop) and a systematic study of UQ evaluation metrics (Pattern Recognition Letters, 2024). Ongoing projects extend this line to flow-based pose modeling and collision-aware planning in autonomous driving.

\textbf{Research experience aligned with the internship workflow.}
I am comfortable with the full ML research pipeline your description outlines: identifying research gaps, reviewing state-of-the-art methods, prototyping novel approaches, running experiments at scale, and writing for top-tier venues. I develop primarily in \textbf{PyTorch} and \textbf{JAX}, and I have experience designing experimental protocols and evaluation metrics---skills I also apply as a peer reviewer for IJCV, CVPR, and AAAI. As primary advisor to nine master's students (2023--2026), I collaborate regularly in a multi-researcher setting and value constructive feedback throughout a project.

\textbf{Why Apple MLR Paris.}
I am drawn to Apple's emphasis on open research with product-scale impact, and to the opportunity to collaborate with MLR colleagues across Europe and the US. The topics you highlight---generative modeling, efficient inference, optimization, and data-centric ML---connect naturally to my current and planned work. I would welcome the chance to contribute new findings through publications and, where appropriate, open-source implementations.

Thank you for considering my application. I would be glad to discuss how my background in uncertainty-aware vision and deep learning research can support your team's goals.

\vspace{1em}
Sincerely,\\
Ziliang Xiong

\end{document}
